% !TEX root = ../main.tex
\section{자료 설명}
\label{ch:appendix-data}

이 부록은 \ref{ch:results}장의 바탕이 되는 자료를 설명해요. 다루는 것은 다음과 같아요. 계약 그래프(점이 모두 계약인 그래프)의 이벤트와 층, 다섯 토큰과 그 가격과 허락(allowance, 토큰을 얼마까지 대신 써도 된다는 허락) 금액, 계약 그래프의 계정 종류, 잔액, 라벨, 그리고 발견 순서 검증의 자세한 내용이에요.

\dissertationSubheading[sec:data-events]{뽑아낸 이벤트와 그래프}

주 분석의 이벤트는 다섯 토큰이 2025년 4월 1일부터 2026년 9월 30일까지 남긴 \texttt{Approval} 로그(금액과 상관없이 모두)와 금액이 0이 아닌 \texttt{Transfer} 로그예요. 이 달들은 두 점수 기간과 두 라벨 기간이 덮는 달이에요. 그 수는 \ref{sec:dataset}절에 있어요. 이 이벤트들은 서로 다른, 순서가 있는 주소 쌍 $73{,}058{,}133$개를 이뤄요. 양 끝에 모두 계약이 있는 행만 그래프에 들어가요(표~\ref{tab:graph-layers}). 점 표에는 두 동결에 걸쳐 계약 $1{,}403{,}959$개가 있어요. 잔액 상한이 있든 없든 어느 한 층의 쌍 끝에 있는 모든 계약, 그리고 코호트 목록과 스펜더 목록의 모든 계약이에요. 두 점수 기간에는 사용자 계정에서 계약으로, 또는 종류를 판별하지 않은 주소로 가는 행이 $152{,}257{,}072$개 더 있어요. 어떤 점수도 이 행들을 읽지 않아요. 발견 순서 검증의 중요도는 이 행들로 계산해요(\ref{sec:discovery-protocol}절). 등록 분석은 2023년 10월 1일부터 2026년 6월 30일까지, 한쪽 끝에 계획의 코호트 계약이 있는 모든 ERC-20 토큰의 이벤트를 뽑았어요. 금액이 0인 것을 포함한 허락 로그 $389{,}866{,}882$개와, 금액이 0이 아닌 송금 로그 $1{,}538{,}404{,}067$개예요.

\dissertationSubheading[sec:data-tokens]{토큰, 가격, 허락 금액}

다섯 토큰을 고른 규칙은 \ref{sec:amount-revision}절에 있어요. 표~\ref{tab:usd-tokens}에는 토큰마다 주소, 소수 자릿수, Chainlink 가격 피드(블록체인에 올라오는 가격 정보), $T_{\mathrm{obs}}$에서 계약 사이의 0보다 큰 마지막 허락(가장 최근 허락)을 잔액 상한 전과 후로 센 수, 그리고 하루 가격의 바탕이 된 자료가 있어요.

% Generated by experiments/scripts/export_latex.py (export_latex_c12.py) -- do not edit by hand.
% Source: experiments/data/0-usd-token-list-and-prices/listed_tokens.csv
% Source: experiments/data/2-processed-tables-and-evaluations/revised-contract-12m/describe/describe.json
% Source: experiments/data/0-usd-token-list-and-prices/chainlink_check.json
\begin{table}[htbp]
\centering
\caption[다섯 토큰, 계약 사이의 허락, 그리고 가격.]{이 분석에 쓴 다섯 토큰이에요. 토큰마다 다음을 적었어요. 먼저 체인링크(Chainlink) 가격 피드(시장 위험 등급)를 적었어요. $T_{\mathrm{obs}}$에서 계약 사이에 있는 0보다 큰 마지막 허락의 수와, 그 가운데 잔액 상한을 씌운 뒤에도 0보다 큰 허락의 수도 적었어요. 가격 시계열에서 두 번째 DefiLlama 키로 채운 날, 전날 값으로 채운 날, 하루짜리 오류로 보고 바꾼 날의 수도 적었어요(기간들이 덮는 2025년 4월 1일부터 2026년 9월 30일까지의 548일 가운데). 마지막으로, 그 날들 가운데 점검 날짜 20개에 블록체인에서 읽은 Chainlink 피드와 DefiLlama 가격 사이 절대 차이의 중앙값과 가장 큰 값을 퍼센트로 적었어요. 종류 판별 기간에 다섯 토큰의 이벤트는 1{,}240{,}275{,}042건이고, 그 가운데 603{,}475{,}784건은 양 끝이 모두 계약이에요.}
\label{tab:usd-tokens}
\fitwidth{%
\begin{tabular}{llrlrrrrrr}
\toprule
토큰 & 주소 & 소수 자릿수 & Chainlink 피드 & 허락 & 상한 뒤 & 2번째 키 & 채움 & 오류 & CL 차이 (\%) \\
\midrule
USDC & \texttt{0xaf88\ldots{}5831} & 6 & USDC / USD (낮음) & 255{,}339 & 108{,}916 & 0 & 0 & 0 & 0.01 / 0.02 \\
USDT0 & \texttt{0xfd08\ldots{}cbb9} & 6 & USDT / USD (낮음) & 36{,}468 & 5{,}648 & 0 & 0 & 0 & 0.01 / 0.15 \\
WETH & \texttt{0x82af\ldots{}bab1} & 18 & ETH / USD (낮음) & 30{,}826 & 1{,}395 & 0 & 0 & 0 & 0.04 / 0.14 \\
WBTC & \texttt{0x2f2a\ldots{}5b0f} & 8 & WBTC / USD (낮음) & 11{,}841 & 386 & 0 & 0 & 0 & 0.05 / 0.24 \\
ARB & \texttt{0x912c\ldots{}6548} & 18 & ARB / USD (낮음) & 4{,}808 & 378 & 0 & 0 & 0 & 0.06 / 0.24 \\
\bottomrule
\end{tabular}
}
\end{table}


다섯 토큰은 저마다 2025년 4월 1일부터 2026년 9월 30일까지의 $548$일 하루하루에, 자기 아비트럼 원(Arbitrum One) 주소로 찾은 DefiLlama 가격이 있어요. 그 서비스의 다른 키에서 가져온 날, 전날 값으로 채운 날, 하루짜리 오류로 보고 고친 날은 하나도 없어요. 토큰마다 이 날들 가운데 $20$일에 체인 위에서 읽은 Chainlink 가격 피드와 비교하면, DefiLlama 가격은 토큰마다 중앙값으로 $0.01\%$에서 $0.06\%$만큼 달라요. 가장 크게 달라도 $0.24\%$예요.

무제한 허락(금액 한도가 없는 허락)은 흔해요. $T_{\mathrm{obs}}$에서 계약 사이의 0보다 큰 마지막 허락 $339{,}282$개 가운데 $223{,}892$개($66.0\%$)는 원래 단위(소수점 조정을 하지 않은 가장 작은 단위)로 $1.15\times10^{77}$ 이상이에요. 이것은 $2^{256}-1$에 가까워요. 그리고 무제한 허락을 포함해 $230{,}313$개는 잔액 상한 전에 값이 상한 $U$(이보다 큰 금액은 이 값으로 자름) 이상이에요(토큰마다의 무제한 허락은 표~\ref{tab:allowance-cap}에 있어요). 등록 분석(미리 등록한 그대로 한 분석)은 모든 토큰을 원래 단위 그대로 두었어요. 그 분석에서는 에고 송금의 $99.52\%$와, $T_{\mathrm{obs}}$ 시점의 마지막 허락 $28{,}069{,}418$개 가운데 $99.78\%$가 원래 단위로 10 이상이었어요. 그 범위에서는 $b=1$일 때 $V$가 1과 $10^{-4}$ 안쪽으로만 차이 나요. 그리고 그 분석에서는 0보다 큰 마지막 허락을 해 둔 주인 $13{,}029{,}306$명 가운데 $1{,}532{,}301$명이 무제한 허락과 그보다 작은 허락을 둘 다 해 두었어요.

\dissertationSubheading[sec:data-contract-graph]{계약 그래프: 계정 종류, 스마트 계정, 잔액}

표~\ref{tab:contract-types}에는 계약 그래프의 바탕이 되는 계정 종류가 있어요. 규칙은 \ref{sub:account-types}절에, 읽기는 부록~\ref{sec:account-type-classification}에 있어요.

% Generated by experiments/scripts/export_latex.py (export_latex_c12.py) -- do not edit by hand.
% Source: experiments/data/2-processed-tables-and-evaluations/revised-contract-12m/describe/describe.json
\begin{table}[htbp]
\centering
\caption[계약 그래프의 계정 종류.]{계약 그래프의 계정 종류예요. 종류 판별 기간에 다섯 토큰의 이벤트에는 0 주소가 아닌 서로 다른 주소가 25{,}537{,}280개 있어요. 그 가운데 17{,}477{,}779개는 트랜잭션(블록체인에 보내는 거래 요청)을 보냈고, 그 사실만으로 사용자 계정이에요. 이 표는 블록 513{,}216{,}506에서 \texttt{eth\_getCode}로 읽은 주소 6{,}230{,}270개의 종류를 판별해요. 여기에는 트랜잭션을 보낸 적이 없고 트랜잭션을 보낸 주소도 0 주소도 아닌 상대가 있는 모든 주소, 코호트의 구성원, 그리고 두 가지 감사 표본이 들어요. 이 집합 밖에서 트랜잭션을 보낸 적이 없는 주소는 상대가 모두 사용자 계정이어서, 계약 사이 화살표의 끝이 될 수 없어요. 스마트 계정은 사용자 계정이에요. Safe는 자기 문턱값과 주인들을 답해요. ERC-4337 계정은 표준 EntryPoint를 답하고, \texttt{validateUserOp} 셀렉터를 가지며 paymaster 셀렉터는 없어요. 프록시가 그 셀렉터를 가리는 계정 192{,}656개는 목록에 오른 코드 계열로 옮겼어요(코드 해시로 37{,}337개, 코드 해시와 구현으로 155{,}319개). 읽은 주소 가운데 524{,}454개는 ERC-1967 프록시, 503{,}516개는 ERC-1167 복제, 45{,}362개는 Safe 프록시예요. 트랜잭션을 보낸 주소까지 합치면, 주소 21{,}877{,}279개가 사용자 계정으로, 1{,}829{,}719개가 계약으로 판별돼요.}
\label{tab:contract-types}
\fitwidth{%
\begin{tabular}{llr}
\toprule
그 블록에서 읽은 코드 & 분류 & 주소 수 \\
\midrule
코드가 없는 외부 소유 계정 & 사용자 계정 & 2{,}909{,}509 \\
EIP-7702 위임이 있는 외부 소유 계정 & 사용자 계정 & 1{,}165{,}025 \\
Safe (문턱값과 주인들을 답함) & 사용자 계정 & 46{,}204 \\
ERC-4337 계정, 읽기 규칙으로 판별 & 사용자 계정 & 87{,}157 \\
ERC-4337 계정, 목록에 오른 코드 계열로 판별 & 사용자 계정 & 192{,}656 \\
그 밖의 모든 코드 & 계약 (그래프의 점) & 1{,}829{,}719 \\
\midrule
읽은 주소 & & 6{,}230{,}270 \\
\bottomrule
\end{tabular}
}
\end{table}


읽기만으로는 프록시가 \texttt{validateUserOp} 셀렉터(함수 식별 값)를 감추는 ERC-4337 계정을 놓쳐요. 코드 계열 목록 두 개가 주소 $192{,}656$개를 더 스마트 계정(계약 코드로 된 사용자 계정: Safe, ERC-4337 계정)으로 판별해요. 이 목록들은 코드만 보고 뽑았어요. 그 목록이 옮기는 주소의 라벨, 코호트 소속, 점수는 읽지 않았어요. 목록을 뽑을 때는 후보 계열을 넣은 구조 파일럿(무게 없이 그래프 모양만 본 예비 점검) 그래프와 넣지 않은 그래프의 수가 알려져 있었어요. 두 경우 모두 프로토콜의 규칙은 같은 공개 방식, 같은 다시 계산 일정, 같은 뜨거운 집합 몫을 정했고, 멈춤 표시도 그대로였어요(\ref{sec:discovery-protocol}절). 첫 번째 목록에는 작은 프록시의 코드 해시 $20$개가 있어요. 이 프록시들은 읽기로 구현 계약을 찾지 못했고, 대표 주소가 표준 EntryPoint를 답해요(주소 $37{,}337$개). 두 번째 목록에는 ERC-1967 프록시의 코드 해시와 구현 계약의 짝 일곱 개가 있어요. 이 프록시들의 구현 계약에는 그 셀렉터가 있어요. 그리고 \texttt{entryPoint()}는 아무것도 답하지 않지만, 값을 읽는 함수 \texttt{entrypoint()}는 표준 EntryPoint를 답해요(주소 $155{,}319$개).

규칙이 찾지 못하는 종류의 스마트 계정은 계약으로 남고, 그래프의 점이 돼요. 두 목록을 쓴 뒤에도, 계약으로 판별한 $115{,}084$개는 살펴본 코드에 \texttt{validateUserOp} 셀렉터가 있고 paymaster 셀렉터는 없어요. 그 가운데 $104{,}187$개는 코드 해시 하나를 함께 써요. 그 코드에서 값을 읽는 함수 \texttt{entrypoint()}는 실행이 되돌려지거나 아예 없어요. 그리고 $8{,}793$개는 구현 계약 하나를 돌리는데, 그 EntryPoint는 표준 EntryPoint가 아니에요. 코드가 Safe도 ERC-4337 계정도 아닌 다른 종류의 스마트 계정처럼 보이는 계약도 계약으로 남아요. 서비스(프로토콜)가 결제, 주문, 사용자마다 만드는 계약도 마찬가지예요. 이런 계약은 어떤 규칙도 골라내지 않아요.

동결 시점에 어떤 주인이 어떤 토큰을 가진 잔액은, 체인의 첫 블록부터 동결까지 그 토큰의 \texttt{Transfer} 로그로 그 주인에게 들어온 양에서 나간 양을 뺀 값이에요. 이 값이 음수이거나 $10^{-6}$ 토큰보다 작으면 0으로 정해요(\ref{sub:balances}절). 잔액은 2025년 4월 1일부터 $T_{\mathrm{obs}}$까지 어떤 토큰을 허락한 적이 있는 토큰과 주인의 짝 $6{,}033{,}263$개에 대해 계산했어요. 점검에서는 이 짝 가운데 $1{,}500$개를 블록 $479{,}089{,}705$, 곧 2026년 6월 30일의 마지막 블록에서 \texttt{balanceOf}와 견주었어요. 토큰마다 $T_{\mathrm{obs}}$에서 다시 계산한 잔액이 0인 짝 $150$개와 0보다 큰 짝 $150$개예요. 그 가운데 $1{,}493$개가 맞았어요($99.5\%$, 기준값은 $98\%$). 나머지 일곱 개는 WETH 네 개와 USDT0 세 개이고, 다시 계산한 잔액이 모두 0으로 세는 값이에요. 이 일곱 개에서는 호출이 아무 값도 돌려주지 않았고, 점검은 이것을 맞지 않은 것으로 세요. 답이 온 짝 가운데 다른 것은 하나도 없어요.

표~\ref{tab:allowance-cap}는 잔액 상한이 계약 사이의 마지막 허락에 미친 영향을 토큰마다 보여 줘요. 0보다 큰 채로 남는 허락은 대부분 허락한 금액보다 적은 잔액으로 세요($T_{\mathrm{obs}}$에서 $116{,}723$개 가운데 $114{,}454$개). 두 동결은 USDC에서 달라요. USDC 허락은 $T_{\mathrm{obs}}$에서 $108{,}916$개가, $t_1$에서 $11{,}381$개가 0보다 큰 채로 남아요. 상한을 씌운 뒤 값이 $U$ 이상인 허락은 $T_{\mathrm{obs}}$에서 $150$개, $t_1$에서 $204$개예요. 그래서 C-PR 허락 층의 상한 $U$가 실제로 자르는 허락은 얼마 안 돼요. 라벨은 그래프와 달리 모든 계정을 세요. 표~\ref{tab:label-split}은 두 라벨 기간의 라벨을 저마다, 라벨이 세는 계정의 갈래에 따라 나눠 보여 줘요(\ref{sec:holdout-results}절).

% Generated by experiments/scripts/export_latex.py (export_latex_c12.py) -- do not edit by hand.
% Source: experiments/data/2-processed-tables-and-evaluations/revised-contract-12m/describe/describe.json
\begin{table}[htbp]
\centering
\caption[계약 사이의 마지막 허락과 잔액 상한.]{계약 사이의 마지막 허락과 잔액 상한을 토큰과 동결마다 보여 줘요. 허락은 점수 기간 안에서 (토큰, 주인, 스펜더) 세 짝의 마지막 \texttt{Approval}이고, 값이 0보다 크면 남겨요. 허락은 동결 시점 주인이 가진 그 토큰의 잔액까지만 세요. 무제한 허락은 원래 단위로 $1.1\times10^{77}$ 이상인 허락이에요. 동결 시점에 주인이 그 토큰을 하나도 갖고 있지 않으면 허락은 0으로 잘려요. 잔액이 0보다 크지만 허락보다 작으면 허락은 잔액으로 줄어들어 잘려요. 0보다 큰 값이 남는 허락만 허락 층에 들어가요. 마지막 열은 상한을 씌운 값을 달러로 더한 것이에요. $T_{\mathrm{obs}}$에서 허락의 65.6\%는 0으로 잘리고 66.0\%는 무제한 허락이에요($t_1$에서는 93.6\%와 67.4\%). $t_1$에서 0보다 큰 상한을 씌운 값의 중앙값은 \$0.071이고, 이것이 상수 $m_A$예요.}
\label{tab:allowance-cap}
\fitwidth{%
\begin{tabular}{lrrrrrr}
\toprule
토큰 & 0보다 큰 허락 & 무제한 & 0으로 잘림 & 잔액으로 잘림 & 상한 뒤 0보다 큼 & 상한을 씌운 값 (USD) \\
\midrule
\multicolumn{7}{l}{\emph{$T_{\mathrm{obs}}$에서}} \\
USDC & 255{,}339 & 167{,}068 & 146{,}423 & 106{,}988 & 108{,}916 & $3.3\times10^{7}$ \\
USDT0 & 36{,}468 & 29{,}609 & 30{,}820 & 5{,}522 & 5{,}648 & $6.5\times10^{6}$ \\
WETH & 30{,}826 & 17{,}608 & 29{,}431 & 1{,}270 & 1{,}395 & $2.0\times10^{7}$ \\
WBTC & 11{,}841 & 5{,}592 & 11{,}455 & 341 & 386 & $7.8\times10^{6}$ \\
ARB & 4{,}808 & 4{,}015 & 4{,}430 & 333 & 378 & $8.4\times10^{5}$ \\
다섯 토큰 전체 & 339{,}282 & 223{,}892 & 222{,}559 & 114{,}454 & 116{,}723 & $6.8\times10^{7}$ \\
\midrule
\multicolumn{7}{l}{\emph{$t_1$에서}} \\
USDC & 223{,}685 & 146{,}928 & 212{,}304 & 9{,}640 & 11{,}381 & $3.3\times10^{7}$ \\
WETH & 36{,}916 & 24{,}198 & 35{,}474 & 1{,}294 & 1{,}442 & $3.2\times10^{7}$ \\
USDT0 & 33{,}673 & 27{,}950 & 27{,}666 & 5{,}851 & 6{,}007 & $7.6\times10^{6}$ \\
WBTC & 9{,}853 & 5{,}409 & 9{,}413 & 390 & 440 & $1.1\times10^{7}$ \\
ARB & 4{,}795 & 3{,}760 & 4{,}442 & 305 & 353 & $6.1\times10^{5}$ \\
다섯 토큰 전체 & 308{,}922 & 208{,}245 & 289{,}299 & 17{,}480 & 19{,}623 & $8.4\times10^{7}$ \\
\bottomrule
\end{tabular}
}
\end{table}


% Generated by experiments/scripts/export_latex.py (export_latex_c12.py) -- do not edit by hand.
% Source: experiments/data/2-processed-tables-and-evaluations/revised-contract-12m/holdout-w0/eval_summary.json
% Source: experiments/data/2-processed-tables-and-evaluations/revised-contract-12m/window-w1/eval_summary.json
\begin{table}[htbp]
\centering
\caption[라벨이 세는 계정의 분류에 따라 나눈 라벨.]{두 라벨 기간의 라벨을, 라벨이 세는 계정의 분류에 따라 나누어 보여 줘요. 라벨마다 스펜더 코호트 전체의 합계, 사용자 계정에서 센 부분과 계약에서 센 부분, 계약의 몫, 그리고 각 부분이 0보다 큰 스펜더의 수를 적었어요. 허락 쌍은 그 주인의 분류로, 송금자는 자기 분류로 세요. 몇몇 송금자는 둘 다 아니에요(토큰을 새로 찍을 때의 0 주소). 그래서 두 부분을 더한 값이 합계보다 그만큼 모자랄 수 있어요. 이렇게 나눈 것은 어떤 규칙에도 들지 않는 설명이에요. 라벨은 12개월 분석의 라벨을 바꾸지 않은 것이고, 라벨이 세는 것은 대부분 사용자 계정이에요. 사용자 계정은 계약 그래프의 어떤 점수도 읽지 않아요.}
\label{tab:label-split}
\fitwidth{%
\begin{tabular}{lrrrrrr}
\toprule
라벨 & 합계 & 사용자 계정에서 & 계약에서 & 계약의 몫 & 사용자 계정 부분이 있는 스펜더 & 계약 부분이 있는 스펜더 \\
\midrule
\multicolumn{7}{l}{\emph{W0 ($n=7{,}864$ 스펜더 계약)}} \\
새 허락 쌍 & 748{,}258 & 676{,}525 & 71{,}733 & 9.6\% & 1{,}322 & 876 \\
새 송금자 & 651{,}474 & 562{,}007 & 89{,}458 & 13.7\% & 1{,}044 & 883 \\
허락 거두기 & 61{,}869 & 60{,}111 & 1{,}758 & 2.8\% & 947 & 287 \\
잔액이 빠져나간 주인 & 570{,}724 & 547{,}427 & 23{,}297 & 4.1\% & 1{,}271 & 926 \\
\midrule
\multicolumn{7}{l}{\emph{W1 ($n=7{,}915$ 스펜더 계약)}} \\
새 허락 쌍 & 702{,}275 & 617{,}088 & 85{,}187 & 12.1\% & 1{,}220 & 641 \\
새 송금자 & 618{,}461 & 509{,}194 & 109{,}258 & 17.7\% & 970 & 855 \\
허락 거두기 & 43{,}625 & 42{,}270 & 1{,}355 & 3.1\% & 761 & 279 \\
잔액이 빠져나간 주인 & 528{,}557 & 503{,}788 & 24{,}769 & 4.7\% & 1{,}139 & 847 \\
\bottomrule
\end{tabular}
}
\end{table}


\dissertationSubheading[sec:discovery-detail]{발견 순서 검증: 우주, 덮는 정도, 민감도, 진단}

이 절에는 \ref{sec:discovery-results}절이 읽는 발견 순서 검증의 표들이 있어요. 표에 나오는 말은 \ref{sec:discovery-protocol}절이 정의해요. 발견 순서 검증은 등록 분석의 결과를 안 뒤에 정했고, 계획과 등록 문서의 어떤 규칙도 정하지 않아요.

표~\ref{tab:discovery-sizes}에는 두 동결에서 계약 그래프, 그 나비넥타이 분해, 두 우주의 크기가 뜨거운 집합의 크기와 함께 있어요. 표~\ref{tab:discovery-coverage}은 각 우주가 두 중요도를 얼마나 담는지 보여 줘요. $T_{\mathrm{obs}}$에서 \texttt{both} 방식의 우주는 값으로 $I_T$의 $96.5\%$와 $I_A$의 $93.0\%$를 담아요. 하지만 각 중요도로 위쪽 $1\%$를 이루는 계약은 $64.2\%$와 $38.1\%$만 담아요.

% Generated by experiments/scripts/export_latex.py (export_latex_c12.py) -- do not edit by hand.
% Source: experiments/data/2-processed-tables-and-evaluations/revised-contract-12m/discovery/discovery_tobs_both.json
% Source: experiments/data/2-processed-tables-and-evaluations/revised-contract-12m/discovery/discovery_tobs_out.json
% Source: experiments/data/2-processed-tables-and-evaluations/revised-contract-12m/discovery/discovery_t1_both.json
% Source: experiments/data/2-processed-tables-and-evaluations/revised-contract-12m/discovery/discovery_t1_out.json
\begin{table}[htbp]
\centering
\caption[계약 그래프와 발견 순서 검증의 우주의 크기.]{두 동결에서 본 계약 그래프의 크기와, 발견 순서 검증의 두 우주(크롤이 다닐 수 있는 계약 전체)의 크기예요. 합집합 그래프에는 두 층 가운데 어느 쪽에든 화살표가 있는, 서로 다른 계약의 순서 있는 쌍이 들어 있어요. 나비넥타이 분해는 합집합 그래프의 것이에요. $t_1$에서 $\mathrm{SCC}\cup\mathrm{OUT}$에는 거대 강연결 성분(SCC)의 약연결 성분(WCC)에 든 계약의 절반이 안 되는 계약만 들어 있어요. 그래서 프로토콜의 규칙에 따라, 두 동결 모두에서 우주가 WCC인 \texttt{both} 방식이 주 방식이고 \texttt{out} 방식은 민감도 분석이에요. 뜨거운 집합 $H_q(I)$에는 우주에서 0보다 큰 중요도 $I$가 가장 큰 계약들이 들어 있어요. 그 크기는 동점을 포함해 우주의 $q=1/100$이에요. 그래서 중요도가 0보다 큰 계약이 적으면 더 작아져요.}
\label{tab:discovery-sizes}
\fitwidth{%
\begin{tabular}{lrr}
\toprule
항목 & $T_{\mathrm{obs}}$ (주) & $t_1$ \\
\midrule
$G$의 계약 (두 동결에 같은 번호) & 1{,}403{,}959 & 1{,}403{,}959 \\
상한을 씌운 허락 쌍 & 115{,}006 & 17{,}940 \\
송금 쌍 & 2{,}580{,}709 & 2{,}269{,}073 \\
합집합 그래프: 쌍이 있는 계약 & 1{,}299{,}438 & 1{,}111{,}554 \\
합집합 그래프: 쌍 & 2{,}575{,}632 & 2{,}269{,}421 \\
\midrule
거대 강연결 성분 (SCC) & 265{,}804 & 222{,}592 \\
IN (SCC에 닿음) & 174{,}239 & 169{,}237 \\
OUT (SCC에서 닿음) & 242{,}833 & 109{,}804 \\
쌍이 있는 다른 계약 & 616{,}562 & 609{,}921 \\
\midrule
$|\mathrm{SCC}\cup\mathrm{OUT}|$, \texttt{out} 방식의 우주 & 508{,}637 & 332{,}396 \\
그 안의 합집합 쌍 & 1{,}546{,}906 & 1{,}283{,}946 \\
$|\mathrm{WCC}|$, \texttt{both} 방식의 우주 (주) & 1{,}169{,}044 & 1{,}015{,}867 \\
그 안의 합집합 쌍 & 2{,}378{,}418 & 2{,}121{,}895 \\
WCC 안의 상한을 씌운 허락 쌍 & 16{,}776 & 17{,}741 \\
\midrule
$|H_q(I_T)|$, \texttt{both} 방식 & 13{,}141 & 12{,}697 \\
$|H_q(I_A)|$, \texttt{both} 방식 & 11{,}662 & 8{,}136 \\
$|H_q(I_T)|$, \texttt{out} 방식 & 5{,}087 & 3{,}324 \\
$|H_q(I_A)|$, \texttt{out} 방식 & 5{,}088 & 3{,}324 \\
\bottomrule
\end{tabular}
}
\end{table}


% Generated by experiments/scripts/export_latex.py (export_latex_c12.py) -- do not edit by hand.
% Source: experiments/data/2-processed-tables-and-evaluations/revised-contract-12m/discovery/discovery_tobs_both.json
% Source: experiments/data/2-processed-tables-and-evaluations/revised-contract-12m/discovery/discovery_tobs_out.json
% Source: experiments/data/2-processed-tables-and-evaluations/revised-contract-12m/discovery/discovery_t1_both.json
% Source: experiments/data/2-processed-tables-and-evaluations/revised-contract-12m/discovery/discovery_t1_out.json
\begin{table}[htbp]
\centering
\caption[계약 그래프와 우주가 두 중요도를 얼마나 덮는지.]{계약 그래프와 발견 순서 검증의 우주들이 두 중요도를 얼마나 덮는지 보여 줘요. $I_T$는 점수 기간에 사용자 계정이 한 계약에 보낸 달러 금액이고, $I_A$는 동결 시점에 사용자 계정이 그 계약에 써도 된다고 허락해 둔 달러 금액이에요. 둘 다 계정마다 $U_{\mathrm{user}}$까지만 세고, $I_A$는 그 계정의 잔액까지만 세요. $T_{\mathrm{obs}}$에서 종류를 판별한 계약 가운데 809{,}296개의 중요도가 0보다 크고, 그 가운데 112{,}210개는 $G$ 밖에 있어요($t_1$에서는 715{,}480개와 109{,}094개). 나비넥타이의 부분은 합집합 그래프의 것이고, 더하면 100\%예요. 사용자 계정은 또 $T_{\mathrm{obs}}$에서 1.68 billion 달러(16억 8,000만 달러), $t_1$에서 1.37 billion 달러(13억 7,000만 달러)를, 상대가 모두 사용자 계정이어서 종류를 판별하지 않은 주소들에 보냈어요. 이 주소들은 그래프의 점이 되지 않아요. 그리고 이 합은 덮는 비율이 어디에도 놓지 못하는 유입의 상한이에요.}
\label{tab:discovery-coverage}
\fitwidth{%
\begin{tabular}{lrrrr}
\toprule
 & \multicolumn{2}{c}{$T_{\mathrm{obs}}$} & \multicolumn{2}{c}{$t_1$} \\
\cmidrule(lr){2-3} \cmidrule(lr){4-5}
중요도의 몫 & $I_T$ & $I_A$ & $I_T$ & $I_A$ \\
\midrule
종류를 판별한 모든 계약의 합 & 16.61 billion 달러 & 309.5 million 달러 & 16.62 billion 달러 & 343.8 million 달러 \\
\midrule
거대 SCC & 71.2\% & 50.2\% & 74.1\% & 49.1\% \\
IN & 9.7\% & 1.7\% & 9.0\% & 2.6\% \\
OUT & 10.1\% & 40.4\% & 6.6\% & 38.2\% \\
쌍이 있는 다른 계약 & 8.1\% & 3.4\% & 9.5\% & 6.3\% \\
합집합 쌍이 없는 $G$의 계약 & 0.4\% & 3.6\% & 0.5\% & 3.3\% \\
$G$ 밖의 계약 & 0.4\% & 0.7\% & 0.3\% & 0.5\% \\
\midrule
WCC가 가진 몫 (\texttt{both} 방식) & 96.5\% & 93.0\% & 97.0\% & 93.8\% \\
$\mathrm{SCC}\cup\mathrm{OUT}$이 가진 몫 (\texttt{out} 방식) & 81.3\% & 90.6\% & 80.7\% & 87.3\% \\
중요도 상위 1\% 계약: WCC 안의 몫 & 64.2\% & 38.1\% & 64.8\% & 34.6\% \\
상위 1\%: $\mathrm{SCC}\cup\mathrm{OUT}$ 안의 몫 & 29.6\% & 26.2\% & 32.5\% & 23.7\% \\
\bottomrule
\end{tabular}
}
\end{table}


표~\ref{tab:discovery-sensitivity}는 판정을 내리는 실행의 설정을 한 번에 하나씩 바꿔요. 어느 행에서도 $D_2$의 구간이 0보다 위에 있지 않아요. $H_q(I_A)$에서 $D_1$의 구간은 두 행에서 0보다 아래에 있고, 다른 두 행에서는 0을 담아요. 표~\ref{tab:discovery-diagnostics}에는 세 가지가 있어요. 각 순서 규칙이 동점으로 정하는 고르기의 몫, 이른 방문 가운데 사용자 계정의 수요가 없는 계약으로 가는 방문의 몫, 그리고 각 순서 규칙이 허브(화살표가 아주 많이 들어오는 계약)에 닿는 방문이에요. 허브의 중앙값은 AWP$'$에서는 방문 $950$번째이고, C-PR$'$에서는 $4{,}088$번째예요. $\mathrm{IB}'_U$에서는 $25{,}565$번째예요.

% Generated by experiments/scripts/export_latex.py (export_latex_c12.py) -- do not edit by hand.
% Source: experiments/data/2-processed-tables-and-evaluations/revised-contract-12m/discovery/discovery_tobs_both.json
\begin{table}[htbp]
\centering
\caption[발견 순서 검증: 판정 대비의 민감도.]{$T_{\mathrm{obs}}$의 주 방식에서 본 발견 순서 검증이에요. 두 판정 대비 $D_1$(C-PR$'$ 빼기 BFS)과 $D_2$(C-PR$'$ 빼기 $\mathrm{IB}'_U$)의 민감도를 보여 줘요. 칸마다 값 아래에, 뜨거운 계약들을 짝지어 다시 뽑은 부트스트랩의 95\% 백분위 구간(다시 뽑기 400번, 시드 42)을 적었어요. 행마다 설정을 하나씩만 바꿔요. 바꾸는 것은 면적을 재는 우주의 몫 $x$, 뜨거운 집합의 몫 $q$, 허브(합집합 입차수가 가장 큰 계약들, 뜨거운 집합에서는 빼지만 그래프에는 남김)를 뺀 뜨거운 집합, 1달러 이상인 사용자 계정의 수로 잰 중요도, 뜨거운 계약 대신 출발 계약을 다시 뽑는 부트스트랩, 그리고 첫 출발 계약들에서 순서 규칙을 네 배 자주 다시 계산하는 일정이에요. 우주에서 $I_A$가 0보다 큰 계약은 11{,}662개뿐이에요. 그래서 $I_A$의 뜨거운 집합은 두 몫 $q$에서 똑같아요. 다른 방식과 다른 동결은 표~\ref{tab:discovery-decision}에 있어요. 이 행들에는 판정이 없어요. 발견 순서 검증은 등록 분석의 결과를 안 뒤에 정했고, 계획과 등록 문서의 어떤 규칙도 정하지 않아요.}
\label{tab:discovery-sensitivity}
\fitwidth{%
\begin{tabular}{lcccc}
\toprule
 & \multicolumn{2}{c}{$H_q(I_T)$에서} & \multicolumn{2}{c}{$H_q(I_A)$에서} \\
\cmidrule(lr){2-3} \cmidrule(lr){4-5}
설정 & $D_1$ & $D_2$ & $D_1$ & $D_2$ \\
\midrule
판정 ($x=10\%$, $q=1\%$) & \begin{tabular}[c]{@{}c@{}}+0.0520\\[-1pt]{\footnotesize [+0.0485, +0.0558]}\end{tabular} & \begin{tabular}[c]{@{}c@{}}-0.0051\\[-1pt]{\footnotesize [-0.0087, -0.0010]}\end{tabular} & \begin{tabular}[c]{@{}c@{}}+0.0061\\[-1pt]{\footnotesize [+0.0014, +0.0111]}\end{tabular} & \begin{tabular}[c]{@{}c@{}}-0.0571\\[-1pt]{\footnotesize [-0.0621, -0.0519]}\end{tabular} \\
\midrule
$x=5\%$ & \begin{tabular}[c]{@{}c@{}}+0.0673\\[-1pt]{\footnotesize [+0.0635, +0.0714]}\end{tabular} & \begin{tabular}[c]{@{}c@{}}+0.0026\\[-1pt]{\footnotesize [-0.0010, +0.0064]}\end{tabular} & \begin{tabular}[c]{@{}c@{}}+0.0472\\[-1pt]{\footnotesize [+0.0429, +0.0521]}\end{tabular} & \begin{tabular}[c]{@{}c@{}}-0.0132\\[-1pt]{\footnotesize [-0.0178, -0.0084]}\end{tabular} \\
$x=20\%$ & \begin{tabular}[c]{@{}c@{}}+0.0199\\[-1pt]{\footnotesize [+0.0169, +0.0230]}\end{tabular} & \begin{tabular}[c]{@{}c@{}}-0.0063\\[-1pt]{\footnotesize [-0.0091, -0.0034]}\end{tabular} & \begin{tabular}[c]{@{}c@{}}-0.0518\\[-1pt]{\footnotesize [-0.0575, -0.0460]}\end{tabular} & \begin{tabular}[c]{@{}c@{}}-0.0966\\[-1pt]{\footnotesize [-0.1023, -0.0905]}\end{tabular} \\
$q=5\%$ & \begin{tabular}[c]{@{}c@{}}+0.0132\\[-1pt]{\footnotesize [+0.0121, +0.0144]}\end{tabular} & \begin{tabular}[c]{@{}c@{}}-0.0067\\[-1pt]{\footnotesize [-0.0078, -0.0057]}\end{tabular} & \begin{tabular}[c]{@{}c@{}}+0.0061\\[-1pt]{\footnotesize [+0.0014, +0.0111]}\end{tabular} & \begin{tabular}[c]{@{}c@{}}-0.0571\\[-1pt]{\footnotesize [-0.0621, -0.0519]}\end{tabular} \\
허브 1{,}170개를 뺌 & \begin{tabular}[c]{@{}c@{}}+0.0306\\[-1pt]{\footnotesize [+0.0274, +0.0341]}\end{tabular} & \begin{tabular}[c]{@{}c@{}}-0.0126\\[-1pt]{\footnotesize [-0.0159, -0.0088]}\end{tabular} & \begin{tabular}[c]{@{}c@{}}-0.0019\\[-1pt]{\footnotesize [-0.0065, +0.0026]}\end{tabular} & \begin{tabular}[c]{@{}c@{}}-0.0607\\[-1pt]{\footnotesize [-0.0654, -0.0553]}\end{tabular} \\
수로 잰 중요도 ($N_T$, $N_A$) & \begin{tabular}[c]{@{}c@{}}+0.0188\\[-1pt]{\footnotesize [+0.0163, +0.0211]}\end{tabular} & \begin{tabular}[c]{@{}c@{}}-0.0282\\[-1pt]{\footnotesize [-0.0308, -0.0254]}\end{tabular} & \begin{tabular}[c]{@{}c@{}}+0.0191\\[-1pt]{\footnotesize [+0.0129, +0.0262]}\end{tabular} & \begin{tabular}[c]{@{}c@{}}-0.0473\\[-1pt]{\footnotesize [-0.0537, -0.0408]}\end{tabular} \\
출발 계약을 다시 뽑는 부트스트랩 & \begin{tabular}[c]{@{}c@{}}+0.0520\\[-1pt]{\footnotesize [+0.0320, +0.0717]}\end{tabular} & \begin{tabular}[c]{@{}c@{}}-0.0051\\[-1pt]{\footnotesize [-0.0249, +0.0179]}\end{tabular} & \begin{tabular}[c]{@{}c@{}}+0.0061\\[-1pt]{\footnotesize [-0.0315, +0.0398]}\end{tabular} & \begin{tabular}[c]{@{}c@{}}-0.0571\\[-1pt]{\footnotesize [-0.0873, -0.0221]}\end{tabular} \\
$\rho/4$ (1/200만큼 늘 때마다 다시 계산, 출발 계약 5개) & \begin{tabular}[c]{@{}c@{}}+0.0381\\[-1pt]{\footnotesize [+0.0347, +0.0418]}\end{tabular} & \begin{tabular}[c]{@{}c@{}}-0.0129\\[-1pt]{\footnotesize [-0.0167, -0.0087]}\end{tabular} & \begin{tabular}[c]{@{}c@{}}-0.0134\\[-1pt]{\footnotesize [-0.0181, -0.0082]}\end{tabular} & \begin{tabular}[c]{@{}c@{}}-0.0713\\[-1pt]{\footnotesize [-0.0765, -0.0658]}\end{tabular} \\
\bottomrule
\end{tabular}
}
\end{table}


% Generated by experiments/scripts/export_latex.py (export_latex_c12.py) -- do not edit by hand.
% Source: experiments/data/2-processed-tables-and-evaluations/revised-contract-12m/discovery/discovery_tobs_both.json
% Source: experiments/data/2-processed-tables-and-evaluations/revised-contract-12m/discovery/discovery_tobs_out.json
\begin{table}[htbp]
\centering
\caption[발견 순서 검증: 동점, 헛 방문, 허브 방문.]{$T_{\mathrm{obs}}$에서 본 발견 순서 검증의 크롤 진단이에요. 동점: 고를 때 같은 점수의 묶음에 계약이 둘 이상 있었던 경우의 몫이에요. 주 방식(\texttt{both})과 \texttt{out} 방식에서 구했어요. 동점은 먼저 찾은 계약을, 그다음에는 주소가 작은 계약을 먼저 골라 풀어요. 헛 방문: 처음 1\%와 처음 10\%의 방문 가운데 $I_T = I_A = 0$인 계약에 간 방문의 몫이에요(주 방식). 허브 방문: 주 우주의 허브 1{,}170개에 대해, 각 허브의 출발 계약별 평균 방문 번호의 중앙값이에요. 가장 큰 허브의 합집합 입차수는 441{,}612예요. 수로 정하는 순서 규칙과 허락 층만 걷는 걷기는 거의 모든 선택을 동점으로 정해요.}
\label{tab:discovery-diagnostics}
\fitwidth{%
\begin{tabular}{lccccr}
\toprule
 & \multicolumn{2}{c}{동점으로 정한 선택} & \multicolumn{2}{c}{헛 방문} & 허브 방문 번호의 \\
\cmidrule(lr){2-3} \cmidrule(lr){4-5}
순서 규칙 & \texttt{both} & \texttt{out} & 처음 1\% & 처음 10\% & 중앙값 \\
\midrule
너비 우선 (BFS) & 0.000 & 0.000 & 0.891 & 0.845 & 73{,}219 \\
무작위 & 0.000 & 0.000 & 0.841 & 0.682 & 272{,}806 \\
\midrule
$\mathrm{IB}'_A$ (허락 역링크) & 1.000 & 1.000 & 0.874 & 0.862 & 59{,}089 \\
$\mathrm{IB}'_T$ (송금 역링크) & 1.000 & 0.999 & 0.678 & 0.888 & 25{,}566 \\
$\mathrm{IB}'_U$ (계약 역링크) & 1.000 & 0.999 & 0.675 & 0.888 & 25{,}565 \\
\midrule
EndorseRank$'$ & 0.999 & 0.998 & 0.874 & 0.863 & 60{,}494 \\
EndorseRank-activity$'$ & 0.991 & 0.997 & 0.692 & 0.857 & 66{,}699 \\
AWP$'$ & 0.252 & 0.005 & 0.522 & 0.835 & 950 \\
\midrule
C-PR$'$ ($\lambda=1$) & 1.000 & 0.999 & 0.874 & 0.863 & 60{,}840 \\
C-PR$'$ ($\lambda=0$) & 0.531 & 0.019 & 0.747 & 0.911 & 5{,}156 \\
C-PR$'$ ($\lambda=0.25$) & 0.534 & 0.019 & 0.727 & 0.908 & 4{,}185 \\
C-PR$'$ ($\lambda=0.5$) & 0.535 & 0.023 & 0.719 & 0.900 & 4{,}088 \\
C-PR$'$ ($\lambda=0.75$) & 0.537 & 0.025 & 0.719 & 0.889 & 4{,}203 \\
S-PR$'$ & 0.531 & 0.019 & 0.743 & 0.910 & 4{,}885 \\
\midrule
$\delta'$ (정규화한 들어오는 무게) & 0.634 & 0.283 & 0.497 & 0.650 & 1{,}404 \\
\bottomrule
\end{tabular}
}
\end{table}


그림~\ref{fig:discovery-curves-out}은 $T_{\mathrm{obs}}$에서 \texttt{out} 방식의 곡선을 보여 줘요. 이 방식에는 판정이 걸려 있지 않아요. 이 방식에서 C-PR$'$은 두 중요도 모두에서 너비 우선 크롤보다, 그리고 $\mathrm{IB}'_U$보다 뜨거운 계약에 빨리 닿아요(표~\ref{tab:discovery-decision}). 하지만 송금 층만 걷는 점수들보다 앞서지는 못해요. $H_q(I_T)$에서는 AWP$'$보다 아래에 있고($-0.0617$, $[-0.0709,-0.0526]$), 자기 송금 층인 C-PR$'$ ($\lambda=0$)보다도 아래에 있어요($-0.0305$, $[-0.0343,-0.0264]$). $H_q(I_A)$에서는 AWP$'$과 비슷하고($-0.0016$, $[-0.0091,+0.0063]$), C-PR$'$ ($\lambda=0$)보다 위에 있어요($+0.0153$, $[+0.0108,+0.0206]$).

\begin{figure}[htbp]
\centering
\includegraphics[width=\textwidth]{\figpath{discovery-curves-out.pdf}}
\caption[{발견 순서 검증: 방문이 나가는 화살표만 드러낼 때 방문한 뜨거운 계약의 몫.}]{$T_{\mathrm{obs}}$에서 \texttt{out} 방식으로 한 발견 순서 검증이에요. 이 방식에는 판정이 걸려 있지 않아요. 가로축은 방문한 계약을 우주에 대한 백분율로 나타낸 값이고, 세로축은 방문한 뜨거운 계약의 몫 $P_{ST}$예요. 왼쪽은 사용자 계정 유입 $I_T$의 뜨거운 집합이고, 오른쪽은 사용자 허락 금액 $I_A$의 뜨거운 집합이에요. 이 방식에서 방문은 나가는 화살표만 드러내요. 그리고 우주는 거대 강연결 성분과 거기서 닿을 수 있는 계약들이에요. 곡선은 너비 우선 순서, 무작위 순서, 역링크 수 $\mathrm{IB}'_U$, AWP$'$, C-PR$'$의 출발 계약 20개에 걸친 평균이에요. 색칠한 띠는 $\mathrm{IB}'_U$와 C-PR$'$의 출발 계약들 사이 가장 작은 값과 가장 큰 값이에요. 점선은 이상적인 순서이고, 세로선은 우주의 $10\%$, 곧 $\mathrm{AUC}_{10\%}$의 범위를 나타내요. 발견 순서 검증은 등록 분석의 결과를 안 뒤에 정했고, 계획과 등록 문서의 어떤 규칙도 정하지 않아요.}
\label{fig:discovery-curves-out}
\end{figure}
